\chapter{Faults and scenarios}
\label{ch:faults}

Faults are objects applied to a running plant, never edits to the physics.
This chapter describes the fault taxonomy in \code{pyplant/faults/model.py},
the scenarios in \code{pyplant/faults/scenarios.py} that turn faults into
labelled data, and how to inject faults into the live stack.

\section{Where a fault enters}

There are three places a fault can enter the loop, and which one it uses is
the most important thing about it.

\begin{figure}[htbp]
\centering
\scalebox{0.8}{%
\begin{tikzpicture}[node distance=8mm and 13mm]
  \node[hot] (plc) {PLC\\LIC101};
  \node[box,right=of plc] (valve) {valve\\FCV101};
  \node[box,right=of valve] (plant) {plant physics\\\code{Params}};
  \node[box,right=of plant] (inst) {instruments\\LT, FT, PT, ZT};
  \draw[flow] (plc) -- node[lbl,above]{OP} (valve);
  \draw[flow] (valve) -- (plant);
  \draw[flow] (plant) -- (inst);
  \draw[flow] (inst.south) -- ++(0,-12mm) -| node[lbl,pos=0.25,below]{measurements} (plc.south);
  \node[bad,above=9mm of valve] (cf) {\textbf{command}\\stiction};
  \node[bad,above=9mm of plant] (pf) {\textbf{parameter}\\leak, fouling, wear};
  \node[bad,above=9mm of inst] (sf) {\textbf{signal}\\bias, drift, stuck, noise};
  \draw[flow,badred] (cf) -- (valve);
  \draw[flow,badred] (pf) -- (plant);
  \draw[flow,badred] (sf) -- (inst);
\end{tikzpicture}}
\caption{The three fault entry points, applied by \code{FaultSet.command()},
\code{FaultSet.params()} and \code{FaultSet.measure()}.}
\label{fig:faultentry}
\end{figure}

\begin{description}
  \item[Parameter faults] change the plant itself: the physics really is
        different afterwards.
  \item[Command faults] sit between the PLC and the valve: the stem does not
        do what it is told.
  \item[Signal faults] corrupt what an instrument reports while the process
        is unchanged.
\end{description}

\begin{theory}[Why the entry point matters]
A level transmitter reading 30\,mm high and a genuine 30\,mm rise produce
identical numbers in \tg{LT101}. They differ only in what the \emph{rest} of
the plant does: a real rise changes \tg{LT102} and the outflow, a bias does
not. A model that has never seen both cannot tell them apart. In a closed
loop it is worse: the controller acts on the biased reading and moves the
\emph{true} level 25\,mm down while the displayed level sits on setpoint.
\end{theory}

\section{Timing}

Every fault has a start $t_s$, an optional end $t_e$ and a ramp $T_r$. Its
strength at time $t$ is
\begin{equation}
  s(t) = \begin{cases}
   0 & t < t_s \text{ or } t \ge t_e,\\
   \min\!\big((t-t_s)/T_r,\;1\big) & T_r > 0,\\
   1 & \text{otherwise.}
  \end{cases}
\end{equation}
The ramp is not cosmetic: an abrupt leak and an incipient one over an hour are
different detection problems, and a dataset containing only step changes
trains a detector that is useless on faults that creep.

\section{The fault catalogue}

\begin{center}
\begin{tabularx}{\textwidth}{@{}lllX@{}}
\toprule
Kind & Entry & Magnitude (default) & Effect \\
\midrule
\code{leak} & parameter & \code{cda} $2\times10^{-5}$\,m$^2$ & adds a Torricelli drain on tank 1: $q_\text{leak}=C_dA\sqrt{2gh_1}$ \\
\code{fouling} & parameter & \code{loss} 0.30 & interconnect $C_dA$ multiplied by $(1-\text{loss}\cdot s)$ \\
\code{pump_wear} & parameter & \code{loss} 0.15 & pump head multiplied by $(1-\text{loss}\cdot s)$ \\
\code{stiction} & command & \code{deadband} 0.06, \code{slip} 0.04 & stem sticks inside $S$, jumps by $J$ \\
\code{sensor_bias} & signal & \code{bias} 25 (units) & reading $+\,\text{bias}\cdot s$ \\
\code{sensor_drift} & signal & \code{rate_per_min} 5 & reading $+\,\text{rate}\cdot(t-t_s)/60$ \\
\code{sensor_stuck} & signal & none & reading frozen at onset value \\
\code{sensor_noise} & signal & \code{sigma} 1.0 & reading $+\,\mathcal N(0,(\sigma s)^2)$, seeded \\
\bottomrule
\end{tabularx}
\end{center}

Signal faults take a \code{tag} (default \tg{LT101}); the corrupted value is
clamped to the tag's span, so a bias cannot read beyond the instrument's
range. Several faults can be active together; the label of a sample is the
sorted list of active fault ids joined by \code{+}, or \code{normal}.

\section{Valve stiction in detail}

Stiction is the most control-relevant fault in the set and the hardest one
for a detector fed only \tg{LT101}, because its signature lives in the phase
relationship between \tg{FCV101_OP} and \tg{ZT101}. pyplant implements the
two-parameter data-driven model (deadband $S$ plus slip jump $J$):

\begin{enumerate}
  \item While stuck, the stem position $y$ does not move.
  \item When the command $u$ moves more than $S$ away from $y$, the stem
        breaks free and lands at $u - \operatorname{sign}(u-y)(S-J)$, i.e.\
        it jumps by $J$ past the point it would have reached smoothly.
  \item While moving, the stem follows the command.
  \item On a reversal of the command direction, it sticks again.
\end{enumerate}

\begin{figure}[htbp]
\centering
\begin{tikzpicture}
\begin{axis}[width=0.5\textwidth,height=4.4cm,xlabel={controller output \tg{FCV101_OP} (\%)},
  ylabel={stem position \tg{ZT101} (\%)},xmin=58,xmax=78,ymin=58,ymax=78,
  grid=major,grid style={linegrey},axis equal image=false]
\addplot[slate!50,dashed,domain=58:78] {x};
\addplot[badred,very thick,->] coordinates {
  (62,62)(68,62)(68,66)(74,72)};
\addplot[badred,very thick,->] coordinates {
  (74,72)(68,72)(68,68)(62,62)};
\node[font=\footnotesize] at (axis cs:65,60.6) {stuck: deadband $S$};
\node[font=\footnotesize,anchor=west] at (axis cs:68.3,64) {slip $J$};
\node[font=\footnotesize] at (axis cs:71,73.5) {stuck on reversal};
\node[font=\footnotesize,anchor=east] at (axis cs:67.7,70) {slip $J$};
\end{axis}
\end{tikzpicture}
\caption{Idealised stiction cycle in the OP--stem plane. A healthy valve lies
on the dashed diagonal; a sticky one traces a parallelogram.}
\label{fig:stiction}
\end{figure}

In a closed loop with integral action the result is a \emph{limit cycle}: the
controller keeps winding until the stem slips, the level overshoots, the
controller winds back the other way, the stem sticks on the reversal, and the
pattern repeats. The parameters are fractions of full stroke and the slip
cannot exceed the deadband.

\section{Instrument noise is not a fault}

Every scenario adds Gaussian noise to the transmitters:

\begin{center}
\begin{tabular}{@{}lccccc@{}}
\toprule
Tag & \tg{LT101} & \tg{LT102} & \tg{FT101} & \tg{PT101} & \tg{ZT101} \\
\midrule
$\sigma$ & 0.4\,mm & 0.4\,mm & 0.03\,L/min & 0.08\,kPa & 0.05\,\% \\
\bottomrule
\end{tabular}
\end{center}

It is applied as \emph{benign} \code{SensorNoise} faults: recorded, never
labelled positive. This is not decoration. A perfectly quiet loop never moves
its valve, so stiction cannot produce a limit cycle at all, and a detector
trained on noiseless signals learns a threshold that collapses the moment it
meets a real transmitter.

\section{Scenarios}

A \code{Scenario} is a reproducible run: a duration, a setpoint programme, a
fault timeline, a seed, plant parameters and a noise table. Running it
produces a trajectory in which every sample is labelled.

\begin{figure}[htbp]
\centering
\scalebox{0.8}{%
\begin{tikzpicture}[node distance=7mm]
  \node[term] (s) {\code{Scenario}};
  \node[box,below=of s] (w) {start at steady state, \code{settle()}\\warm-up up to 1800\,s (not recorded)};
  \node[box,below=of w] (r) {reset clock to $t=0$, attach fresh \code{FaultSet}\\(scenario faults + benign noise)};
  \node[box,below=of r] (run) {run \code{duration} seconds at $\Delta t=0.05$\,s,\\PLC every 0.1\,s, log every 1\,s};
  \node[good,below=of run] (out) {\code{ScenarioResult}: trajectory, labels, events};
  \draw[flow] (s) -- (w);
  \draw[flow] (w) -- (r);
  \draw[flow] (r) -- (run);
  \draw[flow] (run) -- (out);
\end{tikzpicture}}
\caption{\code{run_scenario()}. The warm-up stops a model from learning that
``faults happen at the start''.}
\label{fig:runscenario}
\end{figure}

Fault objects are stateful (stiction carries a stem position), so a fresh copy
is built for every run and the scenario's own fault list is never mutated.

\subsection{The library}
\code{library(duration)} returns one representative scenario per fault mode
plus normal operation, each with onset at 40\% of the duration and the
setpoint programme 300, 340, 280, 300\,mm at 0, 900, 1800, 2700\,s:

\begin{center}
\begin{tabular}{@{}ll@{}}
\toprule
Name & Description \\
\midrule
\code{normal} & setpoint moves only, no fault \\
\code{leak_abrupt} & tank 1 leak, step onset \\
\code{leak_incipient} & tank 1 leak opening over fifteen minutes \\
\code{pump_wear} & impeller wear 18\%, 600\,s ramp \\
\code{fouling} & interconnect restriction 35\%, 1200\,s ramp \\
\code{sensor_bias} & \tg{LT101} reads 25\,mm high \\
\code{sensor_drift} & \tg{LT101} drifting up at 1.5\,mm/min \\
\code{sensor_stuck} & \tg{LT101} frozen \\
\code{stiction} & FCV101 sticking, deadband 6\%, slip 4\% \\
\bottomrule
\end{tabular}
\end{center}

\subsection{The grid}
\code{grid(duration, magnitudes, repeats)} sweeps each fault mode across
\code{magnitudes} scale factors from 0.5 to 1.5, repeats each with different
seeds, and draws a random onset between 25\% and 60\% of the duration. It also
adds \code{repeats} normal runs. A model therefore cannot learn ``the fault
starts at 24 minutes'' instead of learning the fault. The number of scenarios
is
\begin{equation}
  N = \text{repeats}\,\big(1 + 8\times\text{magnitudes}\big),
\end{equation}
so the Standard preset (3 magnitudes, 2 repeats) builds 50 scenarios.
\code{default_dataset()} then tops the normal runs up to \code{normal_runs}
(10 for Standard), giving 58 runs in all.

\section{Injecting faults into the live stack}

The plant service exposes an injection API (Chapter~\ref{ch:api}); the
workbench \ui{Faults} tab is a form over it.

\begin{enumerate}
  \item Start the stack from \ui{Runtime} and wait for all services to be up.
  \item On \ui{Faults}, choose a kind, a magnitude, a ramp and, for sensor
        faults, a tag. Times are in plant seconds from now.
  \item Press \ui{Inject}. The fault appears in the active list with its id.
  \item Watch the HMI: the PV may not move at all (integral action), but the
        OP and the health index will.
  \item Clear one fault or all of them. \code{GET /api/truth} always returns
        the true state, including what the instruments do not show.
\end{enumerate}

From Python:
\begin{lstlisting}
from pyplant.faults.model import Leak
from pyplant.faults.scenarios import Scenario, run_scenario
sc = Scenario("my_leak", duration=1800,
              faults=[Leak("leak", t_start=600, cda=3e-5, ramp=300)])
res = run_scenario(sc)
df = res.frame()        # time series + label + faulty columns
print(res.events)       # ground truth
\end{lstlisting}
